% ============================================================
%  Rozdział 7 — Poszkoleniówkowe PiO
% ============================================================
\section{Poszkoleniówkowe PiO}

Jeśli tutaj dotarłeś/aś, oznacza to, że jesteś pełnoprawnym członkiem klubu i masz takie
same przywileje i obowiązki jak ci długoletni. Regulaminy znajdziesz na naszym Discordzie,
na kanale \texttt{\#przydatne-informacje}.

Wzięcie udziału w szkoleniu gwarantuje opłaconą składkę na najbliższe pół roku.
Co pół roku skarbnik PrzeWrotki przypomina o jej opłaceniu — wynosi ona \textbf{100~zł/semestr}.
Przewidziane są benefity za terminowe płacenie!

% ------------------------------------------------------------
\subsection*{Komunikacja w klubie — Discord}
\addcontentsline{toc}{subsection}{Komunikacja w klubie — Discord}
% ------------------------------------------------------------

Wszystkie wyjazdy klubowe, integracje, pytania i zmiany organizacyjne ogłaszane są
na klubowym Discordzie. Jest on podzielony na kanały — przykładowo, cotygodniowe pływanie
na Odrze jest na \texttt{\#wroclaw-plywanie}, a dalsze wyjazdy w góry — na \texttt{\#whitewater}.

Z początku ilość kanałów może wydawać się przytłaczająca, ale z czasem wszystko stanie się
jasne. Dostęp do Discorda otrzymasz tuż po skończeniu szkoleniówki.

% ------------------------------------------------------------
\subsection*{Godzinki — co to?}
\addcontentsline{toc}{subsection}{Godzinki — co to?}
% ------------------------------------------------------------

Godzinki są naszą walutą klubową. Za godzinki mamy umożliwione wypożyczenie sprzętu
i korzystanie z klubowych atrakcji, a możemy je zarobić uczciwą pracą — przy naprawach
sprzętowych, organizacji spływów, szkoleniu kolejnych pokoleń kajakarzy i wielu innych.

Wszystkie szczegóły i cennik znajdują się w regulaminach na Discordzie.
\textbf{Każdy nowy członek klubu dostaje na swoje konto 6 godzinek.}

% ------------------------------------------------------------
\subsection*{Jak wypożyczyć sprzęt?}
\addcontentsline{toc}{subsection}{Jak wypożyczyć sprzęt?}
% ------------------------------------------------------------

Do zarządzania sprzętem służy nasza autorska aplikacja przeglądarkowa — \textbf{PrzeWrotkApp}:

\begin{center}
  \url{https://app.przewrotka.org/}
\end{center}

Możesz ją dodać do ekranu głównego na swoim telefonie — będzie wtedy działać jak klasyczna
aplikacja. Po skończonej szkoleniówce otrzymasz wiadomość na maila z instrukcjami logowania.

\begin{tipbox}
  Jeśli zauważyłeś (lub spowodowałeś :P) jakiekolwiek wady sprzętu — dodaj je
  w komentarzach i zgłoś to sprzętowcowi, żeby nie nadziać następnej osoby
  na nieprzyjemną niespodziankę!
\end{tipbox}

% ------------------------------------------------------------
\subsection*{Jak zarobić godzinki?}
\addcontentsline{toc}{subsection}{Jak zarobić godzinki?}
% ------------------------------------------------------------

Wiele wydarzeń ogłaszanych na Discordzie wymaga pomocy. Jeśli wykażesz się inicjatywą,
zostaniesz doceniony! Możesz też samodzielnie zorganizować jakieś wydarzenie klubowe
(kajakowe lub nie) — szczegóły w cenniku na \texttt{\#przydatne-informacje}.

Dodatkowo, w zakładce \emph{Zarób} w aplikacji znajdziesz sprzęty z uszkodzeniami
czekające na opiekę.

% ------------------------------------------------------------
\subsection*{Kontakt z zarządem}
\addcontentsline{toc}{subsection}{Kontakt z zarządem}
% ------------------------------------------------------------

\begin{description}[leftmargin=1.5em, font=\bfseries]
  \item[Mateusz (Blue) Soszyński — Prezes]
    Kwestie organizacyjne.
  \item[Kacper (Kacperek) Dąbek — Sprzętowiec]
    Dobór, wypożyczenie i naprawy sprzętu.
  \item[Karolina (Facetka) Sienkiewicz — Szkoleniowiec]
    Dokształcanie się i bycie lepszym kajakarzem.
  \item[Zuzanna Ejsymont — Skarbnik]
    Wszystkie kwestie finansowe.
  \item[Alicja Pernal — Godzinkowa]
    Zgłaszanie wykonanej pracy i rozliczanie godzinek.
  \item[Zuzanna Ejsymont, Zuzanna Kawczyńska, Filip Madzia, Michał Markiewicz (Fantom) — PR]
    Instagram i Facebook. Chętnie przyjmują zdjęcia i opowieści z wypraw.
\end{description}

% ------------------------------------------------------------
\subsection*{Możliwości pływania}
\addcontentsline{toc}{subsection}{Możliwości pływania}
% ------------------------------------------------------------

Bądź na bieżąco z Discordem — każde wydarzenie klubowe ląduje właśnie tam. Poza wyjazdami
ogłaszane są też cykliczne spotkania na stacjonarne ćwiczenie eskimosek na Odrze oraz
spotkania towarzyskie na barce.

Pamiętaj, że Ty też możesz organizować wyjazdy i spotkania — jeśli nie wiesz jak,
rzuć pomysł na Discordzie, chętni się znajdą.

% ------------------------------------------------------------
\subsection*{Gdzie szukać informacji o rzekach?}
\addcontentsline{toc}{subsection}{Gdzie szukać informacji o rzekach?}
% ------------------------------------------------------------

\begin{itemize}[leftmargin=1.5em, itemsep=4pt]
  \item Zapytać bardziej doświadczonych kajakarzy!
  \item Strony internetowe:
    \href{http://wiki.bystrze.pl}{\texttt{wiki.bystrze.pl}},
    \href{http://rivermap.org}{\texttt{rivermap.org}},
    \href{http://wkw.wloclawek.pttk.pl/kajszl.html}{\texttt{wkw.wloclawek.pttk.pl}}
  \item Aplikacje: \emph{Zwałkowe Pegle}, \emph{RiverApp}, \emph{whitewater.guide}
\end{itemize}

% ------------------------------------------------------------
\subsection*{Planowanie spływu jako komandor}
\addcontentsline{toc}{subsection}{Planowanie spływu jako komandor}
% ------------------------------------------------------------

Przydatne narzędzia:

\begin{itemize}[leftmargin=1.5em, itemsep=2pt]
  \item Notatka z zaznaczonymi miejscami charakterystycznymi (mosty, progi, łatwo
    rozpoznawalne punkty)
  \item Pinezki Google Maps
  \item Grupa na Messengerze
  \item Excel z uczestnikami i numerami telefonów
\end{itemize}

% ------------------------------------------------------------
\subsection*{Jak jeszcze mogę wesprzeć klub?}
\addcontentsline{toc}{subsection}{Jak jeszcze mogę wesprzeć klub?}
% ------------------------------------------------------------

Możesz przekazać \textbf{1{,}5\% podatku} na rzecz WKK PrzeWrotka:

\begin{center}
  \begin{tabular}{ll}
    \textbf{KRS:}              & 0000270261 \\
    \textbf{Cel / tytuł:}      & WKK Przerwotka 21794 \\
    \textbf{Numer konta:}      & 35 2030 0045 1110 0000 0382 5640 \\
  \end{tabular}
\end{center}